\documentclass[oneside,12pt,a4paper]{book}

\usepackage[T1]{fontenc} %normes d'encodage
\usepackage[utf8]{inputenc} %normes d'encodage
\usepackage[inner=25mm,outer=25mm,top=30mm,bottom=20mm,headheight=13.6pt]{geometry} %marges
\usepackage{lmodern} %correction d'affichage
\usepackage{graphicx} %insertion d'images
\usepackage{csquotes} %guillemets français
\usepackage{sectsty} %modifier l'apparence des titres
\usepackage{fancyhdr} %en tête et pieds de page personnalisables
\usepackage{etoolbox} %idem
\usepackage{color} %gestion des couleurs
\usepackage{setspace} %interligne
\usepackage{natbib} %pour pouvoir ne citer que l'année
\usepackage[french]{babel} %langue française
\usepackage{chngcntr} %compteurs modifiables
\usepackage{amsmath} %package maths
\usepackage{amssymb} %package maths
\usepackage[lite]{mtpro2} %package maths
\usepackage{url} %gestion des liens hypertextes
\usepackage{blindtext} %génération de texte aléatoire : \Blindtext
\usepackage{tikz} %dessin (pour la page de garde)

\counterwithout{figure}{chapter} %permet un numérotage des figures indépendant des chapitres
\frenchspacing %adaptation normes françaises
\FrenchFootnotes %adaptation normes françaises
\onehalfspacing %Interligne 1,5
\setcounter{secnumdepth}{4} %chapitrage à 4 niveaux (ex: 1.2.2.3. Lorem Ipsum)
\setcounter{tocdepth}{4} % chapitrage à 4 niveaux dans le sommaire

%%%%%%%%%%%%%%%%%%%%%%%%%%%%%%%%%%%%%%%%%%%%%%%%%%%%%%%%%%%%%%%
%%%%%%%%%%%%%%%%%%% Chapitrage personnalisé %%%%%%%%%%%%%%%%%%%
% Inspiré du thème Bredele

\makeatletter
\def\thickhrulefill{\leavevmode \leaders \hrule height 1ex \hfill \kern \z@}
\def\@makechapterhead#1{%
  \vspace*{-30\p@}%
  {\parindent \z@ \raggedleft \reset@font
            \scshape \@chapapp{} \thechapter
        \par\nobreak
        \interlinepenalty\@M
    \Huge \bfseries #1\par\nobreak
    %\vspace*{1\p@}%
    \hrulefill
    \par\nobreak
    \vskip 50\p@
  }}
\def\@makeschapterhead#1{%
 \vspace*{-50\p@}%
  {\parindent \z@ \raggedleft \reset@font
            \scshape \vphantom{\thechapter}
        \par\nobreak
        \interlinepenalty\@M
    \Huge \bfseries #1 \par\nobreak
    %\vspace*{1\p@}%
    \hrulefill
    \par\nobreak
    \vskip 30\p@
  }}

\fancyhf{} 

\appto\mainmatter{\pagestyle{fancy}
\renewcommand{\sectionmark}[1]{\markright{\textit{\thesection.\ #1}}}
\renewcommand{\chaptermark}[1]{\markboth{\textit{#1}}{}}
\fancyhead[L,R]{\small\thepage}
    \fancyhead[R]{\small\rightmark}
    \fancyhead[L]{\small \leftmark}
    \fancyfoot[C]{\thepage}
}
%%%%%%%%%%%%%%%%%%%%%%%%%%%%%%%%%%%%%%%%%%%%%%%%%%%%%%%%%%%%%
%%%%%%%%%%%%%%%%%%%%%%%%%%%%%%%%%%%%%%%%%%%%%%%%%%%%%%%%%%%%%


\title{Thèse Paul-Valery}

\begin{document}


%%%%%%%%%%%%%%%%%%%%%%%%%%%%%%%%%%%%%%%%%%%%%%%%%%%%%%%%%%%%
%%%%%%%%%%%%%%%%%% Page de garde %%%%%%%%%%%%%%%%%%%%%%%%%%%
\thispagestyle{empty}
\begin{tikzpicture}[remember picture, overlay, shift=(current page.south)]
  \draw[gray] (current page.north west) rectangle (-4.5,0) ;
  \shade[top color=white , bottom color=gray] (current page.north west) rectangle (-4.5,0) ;
\end{tikzpicture}
\begin{tikzpicture}[remember picture, overlay, shift=(current page.north west)]
  \node[inner sep=0pt] (collegeDoctoral) at (3,-4)
    {\includegraphics[width=.225\textwidth]{collegeDoc}};
  \node[inner sep=0pt] (PaulValery) at (3,-26)
    {\includegraphics[width=.25\textwidth]{paulVa}};
  \node[inner sep=0pt] (thèse) at (13,-3)
    {\includegraphics[width=0.8\textwidth]{these}};
\node[align=center] at (13.5,-10) {
\\ \large 
Délivrée par \textbf{
Université Paul-Valery %%%% UNIVERSITÉ %%%%
}\\ \\ \\ \large 
Préparée au sein de l'école doctorale
XX %%%% ÉCOLE DOCTORALE %%%%
\\ \large 
Et de l'unité de recherche
XX %%%% UNITÉ DE RECHERCHE %%%%
\\ \\ \large 
Spécialité : \textbf{
XX %%%% SPÉCIALITÉ %%%%
}\\ \\ \\ \\ \\ \large 
Présentée par : \textbf{
XX %%%% NOM %%%%
}\\};
\node[draw, align=center] at (13.5,-17) {\, \hspace{12cm}
\\ \Large \textbf{ 
TITRE DE LA THÈSE} \\ \Large \textbf{ 
SUR DEUX LIGNES   %%%% TITRE %%%%
}\\ \large \textbf{
sous-titre %%%% SOUS-TITRE %%%%
}};
\node[align=center] at (13.5,-20) {Soutenue le
XX %%%% DATE %%%%
devant le jury composé de};
\node[text width=12cm] at (13.5,-23) { %%%% JURY %%%%
Civilité Prénom NOM, Grade, Etablissement \hfill Statut Jury \\
Civilité Prénom NOM, Grade, Etablissement \hfill Statut Jury \\
Civilité Prénom NOM, Grade, Etablissement \hfill Statut Jury \\
Civilité Prénom NOM, Grade, Etablissement \hfill Statut Jury \\};
\end{tikzpicture}
%%%%%%%%%%%%%%%%%%%%%%%%%%%%%%%%%%%%%%%%%%%%%%%%%%%%%%%%%%%%%
%%%%%%%%%%%%%%%%%%%%%%%%%%%%%%%%%%%%%%%%%%%%%%%%%%%%%%%%%%%%%


\newpage

\chapter*{Résumé}
\thispagestyle{empty}

Résumé et mots-clés %%%% INTÉGRER ICI LE RÉSUMÉ %%%%

\chapter*{Remerciements}
\thispagestyle{empty}

Au castor %%%% INTÉGRER ICI LES REMERCIEMENTS %%%%

\renewcommand{\contentsname}{Sommaire}
\renewcommand\thepage{}
\tableofcontents
\thispagestyle{empty}

\mainmatter
\chapter*{Introduction}
\addcontentsline{toc}{chapter}{Introduction}
\chaptermark{Introduction}

%%%%%%%%%%%%%%%%%%%%%%%%%%%%%%%%%%%%%%%%%%%%%%%%%%%%%%%%%%
%%%%%%%%%%%%%%%%%%%%% Manuscrit %%%%%%%%%%%%%%%%%%%%%%%%%%


Ceci est une introduction

\chapter{Titre de chapitre}

intro du premier chapitre

\section{Titre de section}

texte de la section

\begin{figure*}[ht]
\centering\includegraphics[width=0.5\linewidth]{paulVa}
\caption{légende de la figure}
\end{figure*}


exemple de citation de \cite{Sbirre1906}

\bigskip

Lorem ipsum dolor sit amet, consectetuer adipiscing elit. Maecenas fermentum, elit non lobortis cursus, orci velit suscipit est, id mollis turpis mi eget orci. Ut aliquam sollicitudin metus. Mauris at sapien sed sapien congue iaculis. Nulla lorem urna, bibendum id, laoreet iaculis, nonummy eget, massa \citep{Pierre1901}. Phasellus ullamcorper commodo velit. Class aptent taciti sociosqu ad litora torquent per conubia nostra, per inceptos hymenaeos. Phasellus est. Maecenas felis augue, gravida quis, porta adipiscing, iaculis vitae, felis. Nullam ipsum. Nulla a sem ac leo fringilla mattis. Phasellus egestas augue in sem. Etiam ac enim non mauris ullamcorper scelerisque. In wisi leo, malesuada vulputate, tempor sit amet, facilisis vel, velit. Mauris massa est, sodales placerat, luctus id, hendrerit a, urna. Nullam eleifend pede eget odio. Duis non erat. Nullam pellentesque.

Lorem ipsum dolor sit amet, consectetuer adipiscing elit. Maecenas fermentum, elit non lobortis cursus, orci velit suscipit est, id mollis turpis mi eget orci. Ut aliquam sollicitudin metus. Mauris at sapien sed sapien congue iaculis. Nulla lorem urna, bibendum id, laoreet iaculis, nonummy eget, massa. Phasellus ullamcorper commodo velit. Class aptent taciti sociosqu ad litora torquent per conubia nostra, per inceptos hymenaeos. Phasellus est. Maecenas felis augue, gravida quis, porta adipiscing, iaculis vitae, felis. Nullam ipsum. Nulla a sem ac leo fringilla mattis. Phasellus egestas augue in sem. Etiam ac enim non mauris ullamcorper scelerisque. In wisi leo, malesuada vulputate, tempor sit amet, facilisis vel, velit. Mauris massa est, sodales placerat, luctus id, hendrerit a, urna. Nullam eleifend pede eget odio. Duis non erat. Nullam pellentesque.

Lorem ipsum dolor sit amet, consectetuer adipiscing elit. Maecenas fermentum, elit non lobortis cursus, orci velit suscipit est, id mollis turpis mi eget orci. Ut aliquam sollicitudin metus. Mauris at sapien sed sapien congue iaculis. Nulla lorem urna, bibendum id, laoreet iaculis, nonummy eget, massa. Phasellus ullamcorper commodo velit. Class aptent taciti sociosqu ad litora torquent per conubia nostra, per inceptos hymenaeos. Phasellus est. Maecenas felis augue, gravida quis, porta adipiscing, iaculis vitae, felis. Nullam ipsum. Nulla a sem ac leo fringilla mattis. Phasellus egestas augue in sem. Etiam ac enim non mauris ullamcorper scelerisque. In wisi leo, malesuada vulputate, tempor sit amet, facilisis vel, velit. Mauris massa est, sodales placerat, luctus id, hendrerit a, urna. Nullam eleifend pede eget odio. Duis non erat. Nullam pellentesque.

\chapter{Autre titre de chapitre}



\section{Autre titre de section}
\subsection{Titre de sous-partie}
\section{Encore un titre original}
\cite{Simon1903} et \cite{Kohler1903}, pour ne citer qu'eux.

\newpage
\addcontentsline{toc}{chapter}{Bibliographie}
\pagestyle{plain}
\bibliographystyle{apalike} % normes APA6
\bibliography{biblio.bib} % fichier bibtex
\thispagestyle{plain}
\pagestyle{fancy}

\appendix
\chapter{Première annexe}
\pagenumbering{Roman}
\setcounter{page}{1}

Ceci est en annexe

\end{document}
